\documentclass[11pt]{article}
\usepackage[margin=1.0in]{geometry}
\usepackage{booktabs}
\usepackage{amsmath}
\usepackage[T1]{fontenc}
\usepackage{lmodern}
\usepackage{microtype}
\usepackage{graphicx}
\usepackage[hidelinks]{hyperref}
\usepackage{underscore}

\title{The curated \texttt{p\_mix} sweep\\
\large Eight arms, what they show, and where the working size should sit}
\date{31 July 2026}
\author{}

\begin{document}
\maketitle

\begin{abstract}
\noindent Eight \texttt{fmix} arms on the same 100k-gate circuit, varying only
\texttt{p\_mix} from 0.05 to 0.40, with the curated DB and \texttt{--db-advance}
both on. The sweep was run to choose an operating point for long runs. Its main
results: transport and positional reach are \emph{flat} across a wide band of
working sizes while \texttt{dmin} rises steeply, giving a clear optimum near
$k \approx 2$; \texttt{p\_mix} mostly buys \emph{position along a common
trajectory} rather than a different quality of mixing; \texttt{dmin} is the sole
exception and has a threshold near $\texttt{p\_mix} \approx 0.10$; and the
affine reconstruction ridge is unmoved by any of it. A later selection-entropy
instrument --- the bits injected through \emph{which} gates were spliced ---
independently confirms the same working size, and shows that fewer than 28\% of
splices have any choice at all.
\end{abstract}

\section{What was run}

Eight arms, identical except for \texttt{p\_mix} $\in \{0.05, 0.10, 0.15, 0.20,
0.25, 0.30, 0.35, 0.40\}$.

\begin{center}
\begin{tabular}{ll}
\toprule
input & \texttt{pre2\_100k.mpmct1} --- gates 100{,}001--200{,}000 of a fresh
        n=128 Gray gadget \\
      & (100{,}000 gates, 512 wires; the \emph{second} slice, since the gadget's \\
      & opening 100k is measurably atypical) \\
budget & 2{,}000{,}000 moves, reports every 100k \\
MIX & $s_{db}=9$, $p_{convex}=1.0$, $p_{mingen}=0.8$ \\
COMP & $s_{db}=12$, $p_{convex}=0.5$, $p_{mingen}=0$ \\
other & \texttt{p\_db}=1 (thermostat never fires), no twists, canary off, seed 101 \\
instr. & \texttt{--anc-samples 256} (sampled ancestry; exact is impossible at 100k) \\
new & \texttt{--curated}, \texttt{--db-advance} \\
\bottomrule
\end{tabular}
\end{center}

\subsection{The curated contract}

Two things must hold or the results are silently wrong.

\paragraph{Value convention.} The curated store was built with the
pre-\texttt{2ed0222a} \emph{swapped-controls} key convention, so it must be
decoded as such:
\begin{verbatim}
FROZEN_CURATED_VALUE_CONVENTION=legacy-swapped-controls
FROZEN_REGULAR_VALUE_CONVENTION=native
\end{verbatim}
Without this every curated candidate fails verification --- which was the root
cause of the long-standing ``430k candidates, none equivalent'' failure.
Startup must print
\texttt{value conventions: regular=native, curated=legacy-swapped-controls}.
The store identity was verified before use: \texttt{shard\_66} $=6{,}350{,}369$~B,
total $1{,}720{,}215{,}364$~B (the bounded build, not the old 6.49~GB one).

\paragraph{MIX/COMP asymmetry.} This is why higher \texttt{p\_mix} is worth
sweeping at all:
\begin{itemize}
\item \textbf{MIX / expansion} probes the \textbf{curated store first} (forward
key only), taking a random no-larger candidate, else a random minimal one, and
falls back to the regular store only on a class miss.
\item \textbf{COMP / compression is regular-only} --- curated never participates.
\end{itemize}
So curated turns the \emph{expansion} channel into a compressing one. Adding MIX
no longer costs size the way it did without curated.

\texttt{FROZEN\_FILTER} was deliberately left off: it is trajectory-identical
(filters only accelerate misses) and costs ${\sim}25.5$~GB per process, which is
impossible across eight arms.

\section{Results}

\subsection{The two budgets disagree}

\begin{center}
\includegraphics[width=\textwidth]{../reports/ancestry_20260728/curated_pmix_sweep_20260731.png}
\end{center}

All eight ran exactly 2M moves, so the ``equal compute'' comparison is exact.
Because higher \texttt{p\_mix} inflates the circuit, the arms reach very
different \emph{per-gate} work in that budget (17.05 down to 8.02), so the
work-normalised view is given alongside.

\begin{center}
\small
\begin{tabular}{lrrrrrrrr}
\toprule
\texttt{p\_mix} & size & $k$ & eff & \texttt{anc} & incid.\ ($10^9$) & \texttt{cov} & \texttt{ent} & \texttt{dmin} \\
\midrule
0.05 & 128{,}355 & 1.28 & \textbf{17.05} & 44{,}302 & 5.69 & \textbf{0.519} & \textbf{0.799} & 0.038 \\
0.10 & 150{,}699 & 1.51 & 14.70 & 38{,}844 & \textbf{5.85} & 0.461 & 0.771 & 0.070 \\
0.15 & 176{,}429 & 1.76 & 13.01 & 32{,}181 & 5.68 & 0.381 & 0.726 & 0.098 \\
0.20 & 202{,}674 & 2.03 & 11.71 & 26{,}141 & 5.30 & 0.313 & 0.683 & 0.124 \\
0.25 & 235{,}577 & 2.36 & 10.61 & 22{,}119 & 5.21 & 0.268 & 0.643 & 0.147 \\
0.30 & 277{,}404 & 2.77 &  9.57 & 17{,}719 & 4.92 & 0.218 & 0.595 & 0.168 \\
0.35 & 317{,}813 & 3.18 &  8.76 & 14{,}896 & 4.73 & 0.187 & 0.553 & 0.186 \\
0.40 & 369{,}987 & 3.70 &  8.02 & 12{,}220 & 4.52 & 0.158 & 0.507 & \textbf{0.202} \\
\bottomrule
\end{tabular}
\end{center}

At \textbf{matched per-gate work} (eff $=8.02$) the ordering inverts for
transport: incidence rises monotonically $0.74 \to 4.52$, while \texttt{anc},
\texttt{cov} and \texttt{ent} peak at \texttt{p\_mix} $0.15$--$0.25$.

The curated share of successful splices is \textbf{flat at 0.434--0.446} across
the entire sweep. \texttt{p\_mix} does not change how often curated is
\emph{usable}; only how often expansion is attempted.

\subsection{Progression: \texttt{p\_mix} buys position, not quality}

\begin{center}
\includegraphics[width=\textwidth]{../reports/ancestry_20260728/curated_progression_20260731.png}
\end{center}

Panels (b) and (f) are the important ones. Plotted against effective work,
\texttt{anc} and \emph{absolute} positional reach ($\texttt{cov} \times
\text{size}$, i.e.\ gates reached rather than the size-relative fraction) very
nearly \textbf{collapse onto a single curve} regardless of \texttt{p\_mix}. The
arms differ mainly in \emph{how far along that common curve they get}, not in the
shape of the curve. Much of the apparent spread in the endpoint tables is
therefore ``who got further'', not ``who mixed better''.

\paragraph{\texttt{dmin} is the exception, and it has a threshold.}
Panel (d) does not collapse. The arms sit on separated levels, and the
\emph{shapes} differ qualitatively:
\begin{itemize}
\item \texttt{p\_mix} $=0.05$: \texttt{dmin} \textbf{decays} through the run
      ($0.053 \to 0.038$);
\item \texttt{p\_mix} $=0.10$: roughly flat;
\item \texttt{p\_mix} $\ge 0.15$: \texttt{dmin} \textbf{rises} through the run.
\end{itemize}
So there is a threshold near \texttt{p\_mix} $\approx 0.10$ separating
\emph{spending} spelling diversity from \emph{accumulating} it. This is the only
measure on which \texttt{p\_mix} changes the behaviour rather than the rate, and
it is a genuine change from the regular-only runs, where \texttt{dmin} decayed
toward the \texttt{fcompress}-minimal form at every schedule.

\subsection{The reconstruction ridge does not move}

\begin{center}
\includegraphics[width=\textwidth]{../reports/ancestry_20260728/curated_ridge_20260731.png}
\end{center}

Affine (degree-1) reconstruction plates of each final circuit against the
slice's own input --- i.e.\ ``can an adversary holding the input align it to the
mixed circuit''.

\begin{center}
\small
\begin{tabular}{lrrrrrrrr}
\toprule
\texttt{p\_mix} & 0.05 & 0.10 & 0.15 & 0.20 & 0.25 & 0.30 & 0.35 & 0.40 \\
\midrule
depth    & 0.2073 & 0.2070 & 0.2046 & 0.2054 & 0.2030 & 0.2043 & 0.2064 & 0.2000 \\
contrast & 2.445 & 2.475 & 2.460 & 2.422 & 2.434 & 2.388 & 2.441 & 2.381 \\
$\rho$   & 1.000 & 1.000 & 1.000 & 1.000 & 1.000 & 1.000 & 1.000 & 1.000 \\
\bottomrule
\end{tabular}
\end{center}

An $8\times$ range in \texttt{p\_mix} and a $3.7\times$ range in final size move
depth by 3.5\% with no monotone trend, and $\rho$ is exactly $1.000$ everywhere.
The only monotone signal is a 3\% drift in mean $H$ ($0.2560 \to 0.2630$) with a
matching fall in its spread --- the plate flattening very slightly overall,
without touching the ridge.

Combined with the separate finding that depth is unchanged between eff $16$ and
eff $50$ on the earlier long runs, the ridge is so far \textbf{invariant to
everything variable inside DB mixing}: schedule, curated, \texttt{db-advance},
run length and circuit size.

\subsection{\texttt{dmin}'s rise is numerator-driven}

\begin{center}
\includegraphics[width=0.92\textwidth]{../reports/ancestry_20260728/dmin_components_20260731.png}
\end{center}

\texttt{dmin} is a ratio, so a rising ratio could in principle be a shrinking
sample. It is not. Splitting it at 2M moves:

\begin{center}
\begin{tabular}{lrrr}
\toprule
 & \texttt{p\_mix}\,0.05 & \texttt{p\_mix}\,0.40 & span \\
\midrule
denominator \texttt{dmin\_windows} (store held any equivalent) & 1{,}006{,}347 & 1{,}383{,}431 & $1.37\times$ \\
numerator \texttt{dmin\_shorter} (a strictly shorter spelling) &    38{,}199 &   279{,}552 & $7.32\times$ \\
\texttt{dmin} & 0.038 & 0.202 & $5.32\times$ \\
\bottomrule
\end{tabular}
\end{center}

The denominator moves only $1.37\times$ and moves \emph{upward} --- the direction
that would suppress the ratio. So the $5.3\times$ rise is genuinely because
$7.3\times$ more windows admit a shorter spelling. The denominator is also
informative on its own: it rising with \texttt{p\_mix} means MIX-heavy schedules
produce material the store can answer more often, consistent with curated
expansion emitting store-shaped (g57) circuits.

\section{Selection entropy}

A later instrument (\texttt{choice n= multi= mean= bits/splice=}) measures what
\texttt{dmin} could only proxy. \texttt{match\_count} counts everything the store
returns, but what actually injects entropy is the \textbf{eligible} set the
mode's size rule admits, since the winner is drawn uniformly from it. Per
successful splice the mixer now accumulates how many splices had any choice, the
fraction with more than one, the mean branching factor, and the summed $\log_2$
--- \textbf{the bits that entered through \emph{which} gates were spliced, as
distinct from \emph{where} they were spliced}.

The eight arms were re-run with it; final sizes reproduced the originals exactly,
so these are the same trajectories with entropy logged throughout.

\begin{center}
\includegraphics[width=\textwidth]{../reports/ancestry_20260728/selection_entropy_20260731.png}
\end{center}

\begin{center}
\small
\begin{tabular}{lrrrrrr}
\toprule
\texttt{p\_mix} & splices & \texttt{multi} & \texttt{mean} & bits/splice & total bits & bits/gate \\
\midrule
0.05 &   498{,}682 & 0.079 & 1.38 & 0.173 &  86k & 0.672 \\
0.10 &   632{,}151 & 0.115 & 1.55 & 0.251 & 159k & 1.053 \\
0.15 &   724{,}153 & 0.148 & 1.70 & 0.323 & 234k & 1.326 \\
0.20 &   807{,}330 & 0.176 & 1.84 & 0.384 & 310k & 1.530 \\
0.25 &   881{,}629 & 0.204 & 1.99 & 0.448 & 395k & 1.677 \\
0.30 &   947{,}907 & 0.229 & 2.13 & 0.506 & 480k & 1.729 \\
0.35 & 1{,}008{,}625 & 0.252 & 2.26 & 0.561 & 566k & \textbf{1.780} \\
0.40 & 1{,}066{,}045 & 0.275 & 2.39 & 0.615 & 656k & 1.772 \\
\bottomrule
\end{tabular}
\end{center}

\paragraph{Most splices have no choice at all.} Even at \texttt{p\_mix} $=0.40$
only \textbf{27.5\%} of successful splices had more than one eligible candidate,
with a mean branching factor of $2.39$; at $0.05$ it is $7.9\%$ and $1.38$. The
store usually returns a single admissible answer, so the process is considerably
more deterministic than the splice counts suggest.

\paragraph{Entropy injection is front-loaded.} Panel (d): \texttt{bits/splice}
\emph{decays} through every run (${\sim}1.0 \to 0.2$ at the low end) as choice
narrows with approach to local minimality. A long run at fixed \texttt{p\_mix}
injects progressively less entropy per splice.

\paragraph{It confirms the working size independently.} Entropy \emph{density}
(bits per gate) saturates --- 94\% of peak by $k=2.4$, 97\% by $2.8$, and it
\emph{declines} from $k=3.2 \to 3.7$:

\begin{center}
\begin{tabular}{lrr}
\toprule
step in $k$ & bits/gate & transport \\
\midrule
$1.76 \to 2.03$ & $+15.4\%$ & $-6.7\%$ \\
$2.03 \to 2.36$ & $+9.6\%$  & $-1.7\%$ \\
$2.36 \to 2.77$ & $+3.1\%$  & $-5.7\%$ \\
$3.18 \to 3.70$ & $-0.5\%$  & $-4.5\%$ \\
\bottomrule
\end{tabular}
\end{center}

At $k \approx 2.4$ the circuit carries \textbf{94\% of peak entropy density at
89\% of peak transport}, and every further step costs more transport than the
entropy it returns. This is the same knee the transport/\texttt{dmin} analysis
gave, from an independent measure.

\paragraph{Caveat: it is not new information about \texttt{p\_mix}.}
$\mathrm{corr}(\texttt{dmin}, \texttt{bits/splice}) = 0.9987$ across these arms.
Selection entropy \emph{confirms} that \texttt{dmin} was a good proxy rather than
adding an independent signal. Its real value is being denominated in \textbf{bits},
so it can be compared against what state randomization would actually require ---
something \texttt{dmin} could never support.

\section{Conclusions}

\begin{enumerate}
\item \textbf{Curated changes the \texttt{dmin} story.} Without it, \texttt{dmin}
      decayed toward the attacker-computable minimal form under every schedule.
      With it, any \texttt{p\_mix} $\ge 0.15$ accumulates spelling diversity
      instead. That is the clearest qualitative gain in the sweep.
\item \textbf{Most of \texttt{p\_mix}'s apparent effect is effective work.}
      \texttt{anc} and absolute reach follow one common trajectory; the schedule
      sets how far along it a given move budget travels. Comparisons must be made
      at matched per-gate work or they mostly re-measure circuit size.
\item \textbf{Transport is flat over a wide band of working sizes} (Section~4),
      so size can be chosen on other grounds within that band.
\item \textbf{Nothing here moves the ridge.} Structure mixing and state mixing
      are decoupled in this lever set. This is consistent with state
      randomization being the harder pull; it is not yet evidence that it is
      unreachable, but no DB-side lever has moved it at all.
\item \textbf{Choice is scarce.} Fewer than 28\% of successful splices had any
      alternative even at the most MIX-heavy setting, and entropy per splice
      \emph{decays} through every run. Whatever randomness the DB channel
      supplies, it supplies front-loaded and thinly --- which bears directly on
      why the ridge does not move.
\item \textbf{\texttt{cov} must be read in absolute terms across sizes.} Its 64
      buckets are fractions of the current circuit, so a growing circuit shrinks
      \texttt{cov} for reasons unrelated to mixing. Compare
      $\texttt{cov} \times \text{size}$.
\end{enumerate}

\section{Prediction: the steady-state working size}

\paragraph{The execution profile in question.} Expand to $k$ times the initial
size, then set \texttt{p\_mix} so the size holds constant and mix there. The
question is which $k$ to hold at.

\paragraph{Reasoning.} Per fixed compute budget, tabulate what each working size
costs in transport and buys in spelling diversity. Positional reach is given in
absolute gates, since \texttt{cov} alone is size-relative.

\begin{center}
\begin{tabular}{rrrrrr}
\toprule
$k$ & size & incidence & abs.\ reach & \texttt{dmin} & vs.\ $k{=}1.28$ \\
\midrule
1.28 & 128k &  97\% &  96\% & 0.038 & $1.0\times$ \\
1.51 & 151k & \textbf{100\%} & \textbf{100\%} & 0.070 & $1.8\times$ \\
1.76 & 176k &  97\% &  97\% & 0.098 & $2.6\times$ \\
\textbf{2.03} & \textbf{203k} & \textbf{91\%} & \textbf{91\%} & \textbf{0.124} & $\mathbf{3.3\times}$ \\
2.36 & 236k &  89\% &  91\% & 0.147 & $3.9\times$ \\
2.77 & 277k &  84\% &  87\% & 0.168 & $4.4\times$ \\
3.18 & 318k &  81\% &  86\% & 0.186 & $4.9\times$ \\
3.70 & 370k &  77\% &  84\% & 0.202 & $5.3\times$ \\
\bottomrule
\end{tabular}
\end{center}

Transport and absolute reach are \textbf{89--100\% of peak for all $k \le 2.4$},
while \texttt{dmin} climbs $3.9\times$ over the same range. The marginal view
locates the knee precisely:

\begin{center}
\begin{tabular}{lrrr}
\toprule
step in $k$ & incidence & abs.\ reach & \texttt{dmin} \\
\midrule
$1.28 \to 1.51$ & $+2.9\%$ & $+4.3\%$ & $+84.2\%$ \\
$1.51 \to 1.76$ & $-3.0\%$ & $-3.2\%$ & $+40.0\%$ \\
$1.76 \to 2.03$ & $-6.7\%$ & $-5.6\%$ & $+26.5\%$ \\
$2.03 \to 2.36$ & $-1.7\%$ & $-0.5\%$ & $+18.5\%$ \\
\midrule
$2.36 \to 2.77$ & $-5.7\%$ & $-4.2\%$ & $+14.3\%$ \\
$2.77 \to 3.18$ & $-3.7\%$ & $-1.7\%$ & $+10.7\%$ \\
$3.18 \to 3.70$ & $-4.5\%$ & $-1.6\%$ & $+8.6\%$ \\
\bottomrule
\end{tabular}
\end{center}

Every step up to $k=2.4$ returns far more \texttt{dmin} than it costs in
transport. From $k=2.4$ onward the exchange rate turns: the steps cost $4$--$6\%$
of transport for $9$--$14\%$ of \texttt{dmin}, and the \texttt{dmin} gains are
themselves decelerating.

\paragraph{Prediction.}
\begin{center}
\begin{tabular}{ll}
\toprule
working size & $k \approx 2$, usable band $1.8 \le k \le 2.4$ \\
on the 100k slice & \textbf{200{,}000--235{,}000 gates} \\
on the 695k gadget & \textbf{1.4--1.7 million gates} \\
expansion \texttt{p\_mix} & 0.15--0.25 (also the band where \texttt{dmin} trends up) \\
expected \texttt{dmin} at hold & 0.124--0.147 \\
transport cost vs.\ optimum & $\le 11\%$ \\
\bottomrule
\end{tabular}
\end{center}

\textbf{Three} independent arguments land in the same place, which is the main
reason for confidence in the band: the transport/\texttt{dmin} exchange rate
turns at $k \approx 2.4$; \texttt{p\_mix} must be $\ge 0.15$ for \texttt{dmin}
to trend upward at all --- and \texttt{p\_mix} $0.15$--$0.25$ is exactly what
produces $k \approx 1.8$--$2.4$; and (Section~3) selection-entropy density
saturates at 94\% of peak by $k = 2.4$, with every further step costing more
transport than the entropy it returns.

\section{What is not yet measured}

\begin{itemize}
\item \textbf{The hold-phase \texttt{p\_mix} is unknown.} Every arm here grew ---
      even \texttt{p\_mix} $=0.05$ reached $1.28\times$. The size-neutral value
      with curated is therefore \emph{below} 0.05 and has never been measured.
      (Without curated it was ${\approx}0.015$--$0.02$, but curated makes
      expansion far more productive, so that will not transfer.) Three short arms
      at \texttt{p\_mix} $0.01/0.02/0.03$, resumed from a $k\approx2$ state,
      would pin it in under an hour.
\item \textbf{\texttt{dmin} behaviour during the hold is the real unknown.}
      Everything above measures \texttt{dmin} accumulated \emph{while expanding}.
      In the regular-only long runs, holding at low \texttt{p\_mix} made
      \texttt{dmin} decay and flatten (to $0.011$--$0.025$). If curated does the
      same, the banked $0.124$ erodes and the hold policy matters more than the
      working size. If curated sustains it --- plausible, since curated expansion
      injects non-minimal spellings --- then $k \approx 2$ is a genuinely good
      operating point. This is the single measurement to make before committing
      to long runs.
\item \textbf{Everything here is at eff $8$--$17$}, well short of the $50+$ the
      earlier long runs reached, and \texttt{anc} is at most 44\% of the input.
      The orderings could still move.
\item \textbf{Degree-2 plates} were not run on these circuits; the ridge result
      is degree-1 only.
\end{itemize}

\end{document}
